% s-340-p340-m002; original PDF 340
\Aphorism{s-340-p340-m002}{390}
Die ältere Menschheit objectivirt, veräusserlicht, verdinglicht alles Geschehen; der Weg der fortschreitenden Wissenschaft ist der umgekehrte, die Erklärung von innen heraus oder nach innen zu, die Beseitigung der „Erdenreste“, die Reduction äusserer Erfahrungen auf Denkacte und äusserer Erfahrungsgesetze auf Denkbedingungen, kurzum, die Eingliederung der unbekannten Aussenwelt in den Schematismus des einzig Bekannten, des „absoluten Ich“. Der \Sperr{Traum}, der ein Atavismus, ein Zurücksinken in überwundene Entwicklungsstufen ist, stellt die ursprünglichen Denkgewohnheiten wieder her, wie er auch die blind zutappende Schnellmalerei des Causalitätstriebes wieder herstellt, der ohne Wahl und Sichtung die erste beste ihm einfallende Hypothese zu Recht anerkennt, d. h. in die anschauliche Welt hineinconstruirt. So interpretiren wir im Traume auch unsere inneren Zustände als äussere Thatsachen und Objecte: eine Brustbeklemmung als Incubus, eine sexuelle Spannung als erotische Scene, unseren Wunsch weiterzuschlafen als reales Hinderniss des Erwachens — wir treiben Fetischismus und Dämonomagie, wie wir sie in vorgeschichtlichen Zeiten getrieben haben.
\par

